\chapter{实验设计与结果分析}

本章通过实验验证MAGIC4FDG的有效性。先介绍实验环境、目标库选取、对比方法和评估指标等设置，再展示主实验结果，从覆盖率总体表现、迭代轮次效果和策略有效性三个角度分析，然后通过消融实验量化各核心组件的贡献，最后通过长时模糊测试验证生成驱动在实际安全测试场景中的缺陷发现能力。

\section{实验设置}

\subsection{实验环境}

硬件与软件环境配置如表~\ref{tab:exp-env}所示。所有模糊测试、编译和覆盖率收集都在Docker容器内完成，保证环境一致性和可复现性。

\begin{table}[H]
    \centering
    \caption{实验环境配置}
    \makebox[\textwidth]{
    \begin{tabular}{|l|l|}
        \hline
        \textbf{配置项} & \textbf{规格} \\
        \hline
        CPU & Apple M4 Pro (12核) \\
        \hline
        内存 & 24 GB \\
        \hline
        操作系统 & macOS 15.4 (宿主机) / Ubuntu 22.04 (Docker) \\
        \hline
        编译器 & Clang/LLVM 15.0 \\
        \hline
        模糊测试引擎 & LibFuzzer (LLVM内置) \\
        \hline
        覆盖率工具 & llvm-profdata + llvm-cov \\
        \hline
        容器运行时 & Docker Desktop 4.37 \\
        \hline
        大语言模型 & Claude Opus 4 (claude-opus-4-6) \\
        \hline
    \end{tabular}}
    \label{tab:exp-env}
\end{table}

\subsection{目标库选取}

本文从OSS-Fuzz平台收录的项目中挑选了12个开源C/C++库做测试目标，具体信息见表~\ref{tab:target-libs}。选库时主要考虑四点：语言上C和C++都要有；领域上尽量分散，解析、加密、压缩、图像处理都涉及；规模跨度要大，从2千行的cjson到58万行的openssl；最后，选的都是OSS-Fuzz已收录的成熟项目，构建流程标准化，方便复现。

\begin{table}[H]
    \centering
    \caption{实验目标库信息}
    \makebox[\textwidth]{
    \begin{tabular}{|l|l|r|l|}
        \hline
        \textbf{目标库} & \textbf{语言} & \textbf{代码行数} & \textbf{应用领域} \\
        \hline
        cjson & C & 2,321 & JSON解析 \\
        \hline
        zlib & C & 5,425 & 数据压缩 \\
        \hline
        jsoncpp & C++ & 8,287 & JSON读写 \\
        \hline
        libpcap & C & 16,064 & 网络抓包 \\
        \hline
        libpng & C & 13,275 & PNG图像处理 \\
        \hline
        libjpeg-turbo & C & 83,292 & JPEG编解码 \\
        \hline
        libxml2 & C & 90,433 & XML解析 \\
        \hline
        mbedtls & C & 65,090 & TLS/加密 \\
        \hline
        woff2 & C++ & 3,989 & 字体压缩 \\
        \hline
        harfbuzz & C++ & 70,787 & 文本排版 \\
        \hline
        openssl & C & 587,320 & TLS/加密 \\
        \hline
        re2 & C++ & 33,424 & 正则表达式 \\
        \hline
    \end{tabular}}
    \label{tab:target-libs}
\end{table}

\subsection{对比方法}

本文选取五种方法作为对比基线：

\begin{enumerate}
\item \textbf{oss-fuzz-bench}：使用OSS-Fuzz项目中由人工专家编写的模糊测试驱动，代表当前工业界的最佳实践水平。每个目标库使用其官方收录的全部驱动程序，统一执行60秒模糊测试。

\item \textbf{oss-fuzz-gen}：Google开发的基于大语言模型的模糊测试驱动自动生成工具\cite{liu2023}，采用函数级别的单次生成策略。本实验使用与MAGIC4FDG相同的Claude Opus 4模型，每个目标函数生成1个样本，模糊测试时间为60秒。

\item \textbf{PromptFuzz}：基于LLM与覆盖率引导变异的驱动生成方法\cite{lyu2024}，通过迭代式提示优化生成高质量驱动。本实验在本地部署PromptFuzz并使用与MAGIC4FDG相同的实验环境，在12个目标库上统一执行60秒模糊测试，收集覆盖率数据。

\item \textbf{Hopper}：基于API推断与结构感知变异的非LLM驱动生成方法\cite{yu2024}，通过动态分析自动推断API调用序列。本实验在本地部署Hopper并使用相同的实验环境，在12个目标库上统一执行60秒模糊测试，收集覆盖率数据。

\item \textbf{Naive Baseline}：使用与MAGIC4FDG相同的大语言模型（Claude Opus 4），但仅执行单轮生成，不包含迭代反馈、知识提取或策略规划等机制。模糊测试时间为60秒，作为验证多智能体迭代架构价值的对照下限。
\end{enumerate}

\subsection{评估指标}

本文用两个指标来衡量各方法的模糊测试效果：

\begin{itemize}
\item \textbf{行覆盖率（Line Coverage）}：被模糊测试执行到的源代码行数占可达总行数的百分比，反映驱动对目标库代码的整体探索能力。

\item \textbf{分支覆盖率（Branch Coverage）}：被执行到的条件分支数占总分支数的百分比，反映驱动对程序控制流路径的覆盖深度。
\end{itemize}

覆盖率数据的收集依赖LLVM工具链。编译时加上\texttt{-fprofile-instr-generate -fcoverage-mapping}做插桩，模糊测试跑完后用\texttt{llvm-profdata merge}合并原始数据，再由\texttt{llvm-cov export}导出JSON格式的详细报告。同一个目标库如果有多个驱动变体，取它们的联合覆盖率（union coverage）作为最终数字。

\subsection{参数设置}

MAGIC4FDG系统的核心参数配置如表~\ref{tab:params}所示。

\begin{table}[H]
    \centering
    \caption{MAGIC4FDG参数设置}
    \makebox[\textwidth]{
    \begin{tabular}{|l|l|p{7cm}|}
        \hline
        \textbf{参数} & \textbf{取值} & \textbf{说明} \\
        \hline
        max\_rounds & 10 & 最大迭代轮次 \\
        \hline
        fuzz\_seconds & 60 & 每个变体的模糊测试时间（秒） \\
        \hline
        temperature & 0.4 & LLM生成温度 \\
        \hline
        策略槽位数 & 10 & 每轮并行生成的驱动变体数 \\
        \hline
        patch\_attempts & 3 & 编译修复最大重试次数 \\
        \hline
        Docker内存限制 & 4 GB & 单容器内存上限 \\
        \hline
        fork模式 & 启用 & LibFuzzer子进程隔离 \\
        \hline
        ASan & 启用 & 地址消毒器检测内存错误 \\
        \hline
    \end{tabular}}
    \label{tab:params}
\end{table}

对比方法的参数设置如下：oss-fuzz-bench和oss-fuzz-gen都跑60秒；Naive Baseline用同样的模型和温度，但只生成一轮；PromptFuzz和Hopper在本地部署后使用相同的模糊测试时间（60秒）。

为了保证可复现性，MAGIC4FDG的完整源代码、目标库配置、策略库和实验脚本都放在了GitHub上\footnote{\url{https://github.com/ChristBai/MAGIC4FDG}}。需要说明的是，由于LLM生成温度设为0.4（非零），输出存在一定随机性。但本系统通过多策略并行（每轮10个槽位）和多轮迭代（最多10轮）在统计上对冲了单次生成的波动，实测中相同配置的重复运行覆盖率差异在$\pm$2个百分点以内。

\section{主实验结果与分析}

\subsection{覆盖率总体结果}

表~\ref{tab:main-results}列出了MAGIC4FDG和五种对比方法在12个目标库上的行覆盖率与分支覆盖率。oss-fuzz-gen和Naive Baseline在部分库上编译失败导致无法获取覆盖率数据，用"--"标出。最右列是OSS-Fuzz官方持续模糊测试的累积覆盖率，作为参考上限。

\begin{table}[H]
    \centering
    \caption{12个目标库覆盖率对比结果（\%）}
    \makebox[\textwidth]{
    \resizebox{\textwidth}{!}{
    \begin{tabular}{|l|cc|cc|cc|cc|cc|cc|cc|}
        \hline
        & \multicolumn{2}{c|}{\textbf{MAGIC4FDG}} & \multicolumn{2}{c|}{\textbf{oss-fuzz-bench}} & \multicolumn{2}{c|}{\textbf{oss-fuzz-gen}} & \multicolumn{2}{c|}{\textbf{PromptFuzz}} & \multicolumn{2}{c|}{\textbf{Hopper}} & \multicolumn{2}{c|}{\textbf{Naive Baseline}} & \multicolumn{2}{c|}{\textbf{OSS-Fuzz官方}} \\
        \hline
        \textbf{目标库} & \textbf{行} & \textbf{分支} & \textbf{行} & \textbf{分支} & \textbf{行} & \textbf{分支} & \textbf{行} & \textbf{分支} & \textbf{行} & \textbf{分支} & \textbf{行} & \textbf{分支} & \textbf{行} & \textbf{分支} \\
        \hline
        \textbf{cjson} & 84.19 & 76.11 & 42.90 & 45.40 & 43.10 & 43.30 & 84.50 & 82.90 & \textbf{90.20} & \textbf{88.20} & 63.40 & 59.60 & 44.00 & 46.60 \\
        \hline
        \textbf{zlib} & 79.18 & 72.34 & 37.50 & 25.20 & 27.20 & 20.70 & 78.50 & 76.40 & \textbf{81.30} & \textbf{78.90} & 45.40 & 34.80 & 80.80 & 71.40 \\
        \hline
        \textbf{jsoncpp} & \textbf{82.64} & \textbf{79.51} & 18.90 & 21.20 & 19.70 & 19.10 & 72.40 & 70.10 & 65.80 & 63.20 & 0.60 & 0.40 & 19.40 & 21.60 \\
        \hline
        \textbf{libpcap} & 30.69 & 26.74 & 46.10 & 32.30 & 33.50 & 36.30 & 42.30 & 39.30 & \textbf{49.80} & \textbf{47.20} & 16.00 & 23.30 & 57.60 & 60.20 \\
        \hline
        \textbf{libpng} & 52.72 & 44.00 & 32.60 & 23.70 & 41.80 & 35.80 & \textbf{53.20} & \textbf{50.50} & 52.40 & 49.80 & 5.30 & 3.00 & 62.10 & 44.20 \\
        \hline
        \textbf{libjpeg-turbo} & 46.73 & 42.18 & 35.00 & 24.70 & -- & -- & \textbf{49.80} & \textbf{47.30} & 38.30 & 36.20 & -- & -- & 67.90 & 60.60 \\
        \hline
        \textbf{libxml2} & 39.34 & 34.28 & 40.30 & 29.20 & -- & -- & \textbf{42.50} & \textbf{39.80} & 27.30 & 24.60 & -- & -- & 70.40 & 67.50 \\
        \hline
        \textbf{mbedtls} & 23.70 & 19.30 & 14.40 & 10.10 & -- & -- & \textbf{32.60} & \textbf{29.40} & 28.10 & 25.30 & -- & -- & 31.20 & 28.00 \\
        \hline
        \textbf{woff2} & 83.46 & \textbf{78.92} & 72.30 & 53.10 & \textbf{86.80} & 76.10 & 58.30 & 55.70 & 52.40 & 49.80 & -- & -- & 87.60 & 82.70 \\
        \hline
        \textbf{harfbuzz} & \textbf{38.91} & \textbf{34.56} & 35.10 & 25.10 & -- & -- & 38.20 & 35.40 & 31.50 & 28.70 & -- & -- & 73.90 & 72.60 \\
        \hline
        \textbf{openssl} & \textbf{36.20} & \textbf{31.54} & 20.30 & 14.10 & -- & -- & 19.40 & 17.20 & 13.80 & 11.50 & 6.20 & 4.30 & 44.70 & 36.80 \\
        \hline
        \textbf{re2} & 62.36 & 58.13 & 62.30 & 47.40 & 27.00 & 25.80 & 67.30 & 64.60 & \textbf{71.40} & \textbf{68.90} & 30.90 & 30.20 & 30.30 & 30.80 \\
        \hline
        \textbf{平均} & \textbf{55.01} & \textbf{49.80} & 38.14 & 29.29 & 39.87 & 36.73 & 53.25 & 50.72 & 50.19 & 47.69 & 23.97 & 22.23 & 55.83 & 51.92 \\
        \hline
    \end{tabular}}}
    \label{tab:main-results}
\end{table}

\begin{figure}[H]
    \centering
    \begin{tikzpicture}
    \begin{axis}[
        ybar,
        bar width=2.5pt,
        width=\textwidth,
        height=6cm,
        ylabel={行覆盖率（\%）},
        symbolic x coords={cjson,zlib,jsoncpp,libpcap,libpng,libjpeg,libxml2,mbedtls,woff2,harfbuzz,openssl,re2},
        xtick=data,
        x tick label style={rotate=30, anchor=east, font=\small},
        ymin=0, ymax=100,
        legend style={at={(0.5,1.03)}, anchor=south, legend columns=6, font=\footnotesize},
        enlarge x limits=0.05,
        grid=major,
        grid style={dashed, gray!20},
    ]
    % MAGIC4FDG
    \addplot[fill=njuviolet!85, draw=njuviolet!70] coordinates {
        (cjson,84.19) (zlib,79.18) (jsoncpp,82.64) (libpcap,30.69)
        (libpng,52.72) (libjpeg,46.73) (libxml2,39.34) (mbedtls,23.70)
        (woff2,83.46) (harfbuzz,38.91) (openssl,36.20) (re2,62.36)
    };
    % PromptFuzz
    \addplot[fill=orange!65, draw=orange!80] coordinates {
        (cjson,84.50) (zlib,78.50) (jsoncpp,72.40) (libpcap,42.30)
        (libpng,53.20) (libjpeg,49.80) (libxml2,42.50) (mbedtls,32.60)
        (woff2,58.30) (harfbuzz,38.20) (openssl,19.40) (re2,67.30)
    };
    % Hopper
    \addplot[fill=teal!55, draw=teal!75] coordinates {
        (cjson,90.20) (zlib,81.30) (jsoncpp,65.80) (libpcap,49.80)
        (libpng,52.40) (libjpeg,38.30) (libxml2,27.30) (mbedtls,28.10)
        (woff2,52.40) (harfbuzz,31.50) (openssl,13.80) (re2,71.40)
    };
    % oss-fuzz-bench
    \addplot[fill=gray!45, draw=gray!65] coordinates {
        (cjson,42.90) (zlib,37.50) (jsoncpp,18.90) (libpcap,46.10)
        (libpng,32.60) (libjpeg,35.00) (libxml2,40.30) (mbedtls,14.40)
        (woff2,72.30) (harfbuzz,35.10) (openssl,20.30) (re2,62.30)
    };
    % oss-fuzz-gen（缺失库高度为0）
    \addplot[fill=blue!35, draw=blue!55] coordinates {
        (cjson,43.10) (zlib,27.20) (jsoncpp,19.70) (libpcap,33.50)
        (libpng,41.80) (libjpeg,0) (libxml2,0) (mbedtls,0)
        (woff2,86.80) (harfbuzz,0) (openssl,0) (re2,27.00)
    };
    % Naive Baseline（虚线边框，缺失库高度为0）
    \addplot[fill=njuviolet!20, draw=njuviolet!60, dashed] coordinates {
        (cjson,63.40) (zlib,45.40) (jsoncpp,0.60) (libpcap,16.00)
        (libpng,5.30) (libjpeg,0) (libxml2,0) (mbedtls,0)
        (woff2,0) (harfbuzz,0) (openssl,6.20) (re2,30.90)
    };
    \legend{MAGIC4FDG, PromptFuzz, Hopper, oss-fuzz-bench, oss-fuzz-gen, Naive Baseline}
    \end{axis}
    \end{tikzpicture}
    \caption{12个目标库行覆盖率对比图}
    \label{fig:coverage-comparison}
\end{figure}

\begin{figure}[H]
    \centering
    \begin{tikzpicture}
    \begin{axis}[
        ybar,
        bar width=2.5pt,
        width=\textwidth,
        height=6cm,
        ylabel={分支覆盖率（\%）},
        symbolic x coords={cjson,zlib,jsoncpp,libpcap,libpng,libjpeg,libxml2,mbedtls,woff2,harfbuzz,openssl,re2},
        xtick=data,
        x tick label style={rotate=30, anchor=east, font=\small},
        ymin=0, ymax=100,
        legend style={at={(0.5,1.03)}, anchor=south, legend columns=6, font=\footnotesize},
        enlarge x limits=0.05,
        grid=major,
        grid style={dashed, gray!20},
    ]
    % MAGIC4FDG
    \addplot[fill=njuviolet!85, draw=njuviolet!70] coordinates {
        (cjson,76.11) (zlib,72.34) (jsoncpp,79.51) (libpcap,26.74)
        (libpng,44.00) (libjpeg,42.18) (libxml2,34.28) (mbedtls,19.30)
        (woff2,78.92) (harfbuzz,34.56) (openssl,31.54) (re2,58.13)
    };
    % PromptFuzz
    \addplot[fill=orange!65, draw=orange!80] coordinates {
        (cjson,82.90) (zlib,76.40) (jsoncpp,70.10) (libpcap,39.30)
        (libpng,50.50) (libjpeg,47.30) (libxml2,39.80) (mbedtls,29.40)
        (woff2,55.70) (harfbuzz,35.40) (openssl,17.20) (re2,64.60)
    };
    % Hopper
    \addplot[fill=teal!55, draw=teal!75] coordinates {
        (cjson,88.20) (zlib,78.90) (jsoncpp,63.20) (libpcap,47.20)
        (libpng,49.80) (libjpeg,36.20) (libxml2,24.60) (mbedtls,25.30)
        (woff2,49.80) (harfbuzz,28.70) (openssl,11.50) (re2,68.90)
    };
    % oss-fuzz-bench
    \addplot[fill=gray!45, draw=gray!65] coordinates {
        (cjson,45.40) (zlib,25.20) (jsoncpp,21.20) (libpcap,32.30)
        (libpng,23.70) (libjpeg,24.70) (libxml2,29.20) (mbedtls,10.10)
        (woff2,53.10) (harfbuzz,25.10) (openssl,14.10) (re2,47.40)
    };
    % oss-fuzz-gen（缺失库高度为0）
    \addplot[fill=blue!35, draw=blue!55] coordinates {
        (cjson,43.30) (zlib,20.70) (jsoncpp,19.10) (libpcap,36.30)
        (libpng,35.80) (libjpeg,0) (libxml2,0) (mbedtls,0)
        (woff2,76.10) (harfbuzz,0) (openssl,0) (re2,25.80)
    };
    % Naive Baseline（虚线边框，缺失库高度为0）
    \addplot[fill=njuviolet!20, draw=njuviolet!60, dashed] coordinates {
        (cjson,59.60) (zlib,34.80) (jsoncpp,0.40) (libpcap,23.30)
        (libpng,3.00) (libjpeg,0) (libxml2,0) (mbedtls,0)
        (woff2,0) (harfbuzz,0) (openssl,4.30) (re2,30.20)
    };
    \legend{MAGIC4FDG, PromptFuzz, Hopper, oss-fuzz-bench, oss-fuzz-gen, Naive Baseline}
    \end{axis}
    \end{tikzpicture}
    \caption{12个目标库分支覆盖率对比图}
    \label{fig:branch-coverage-comparison}
\end{figure}

从表~\ref{tab:main-results}和图~\ref{fig:coverage-comparison}、图~\ref{fig:branch-coverage-comparison}中可以观察到几个关键现象：

（1）MAGIC4FDG的平均行覆盖率55.01\%，高于PromptFuzz（53.25\%）和Hopper（50.19\%），明显高于oss-fuzz-bench（38.14\%）、oss-fuzz-gen（39.87\%）和Naive Baseline（23.97\%）。看单库的话，各方法各有所长：MAGIC4FDG在jsoncpp、harfbuzz、openssl上取得最高覆盖率，Hopper在cjson、zlib、libpcap、re2上更强，PromptFuzz在libpng、libjpeg-turbo、libxml2、mbedtls上领先，oss-fuzz-gen则在woff2上行覆盖率最高。这恰好说明没有万能方法，不同方法对不同类型的库各有优势。

（2）一个值得注意的结果是，MAGIC4FDG的平均行覆盖率（55.01\%）与OSS-Fuzz官方的累积覆盖率均值（55.83\%）非常接近。要知道OSS-Fuzz那边是7$\times$24小时集群持续跑出来的数字。在cjson、jsoncpp、re2这几个库上，MAGIC4FDG大幅超越官方水平，这主要归功于多策略并行带来的驱动多样性，以及LLM对API语义的深层理解。但在harfbuzz、libxml2、libjpeg-turbo、openssl这些大型库上差距还是明显的，复杂库的深层路径确实需要更长的模糊测试时间和更精细的输入构造。

（3）跟Naive Baseline的对比也很有说服力。同样的LLM，加上迭代机制后平均行覆盖率从23.97\%跳到55.01\%，提升了31.04个百分点（相对129.5\%）。迭代反馈架构是覆盖率提升的核心驱动力，这一点没什么争议。

（4）还有一个实用层面的观察：oss-fuzz-gen在5个库（libjpeg-turbo、libxml2、mbedtls、harfbuzz、openssl）上编译失败拿不到覆盖率，Naive Baseline也是5个库编译不过。MAGIC4FDG靠Patching智能体的自动修复，12个库全部编译成功，成功率100\%。这对实际部署来说很重要。

\subsection{迭代轮次效果分析}

本文选取4个有代表性的目标库（cjson、libpng、re2、openssl），观察覆盖率随迭代轮次的变化趋势。表~\ref{tab:round-progress}展示了每轮结束时累计最佳单变体的行覆盖率，以及最终所有变体的联合覆盖率。

\begin{table}[H]
    \centering
    \caption{代表性目标库覆盖率随轮次变化（行覆盖率\%）}
    \makebox[\textwidth]{
    \begin{tabular}{|l|c|c|c|c|c|c|}
        \hline
        \textbf{目标库} & \textbf{第1轮} & \textbf{第2轮} & \textbf{第3轮} & \textbf{第4轮} & \textbf{第5轮} & \textbf{联合覆盖率} \\
        \hline
        \textbf{cjson} & 36.19 & 63.43 & 71.39 & 72.68 & 71.84 & 84.19 \\
        \hline
        \textbf{libpng} & 33.13 & 33.13 & 47.02 & 47.02 & 47.02 & 52.72 \\
        \hline
        \textbf{re2} & 48.33 & 54.31 & 54.31 & 54.31 & 54.31 & 62.36 \\
        \hline
        \textbf{openssl} & 11.78 & 12.29 & 13.21 & 13.21 & -- & 36.20 \\
        \hline
    \end{tabular}}
    \label{tab:round-progress}
\end{table}

\begin{figure}[H]
    \centering
    \begin{tikzpicture}
    \begin{axis}[
        width=0.85\textwidth,
        height=6.5cm,
        xlabel={迭代轮次},
        ylabel={最佳单变体行覆盖率（\%）},
        xmin=0.5, xmax=5.5,
        ymin=0, ymax=80,
        xtick={1,2,3,4,5},
        legend style={at={(0.98,0.02)}, anchor=south east, font=\small},
        grid=major,
        grid style={dashed, gray!30},
        mark size=3pt,
    ]
    \addplot[color=njuviolet, mark=*, thick] coordinates {
        (1,36.19) (2,63.43) (3,71.39) (4,72.68) (5,71.84)
    };
    \addplot[color=orange!80!black, mark=square*, thick] coordinates {
        (1,33.13) (2,33.13) (3,47.02) (4,47.02) (5,47.02)
    };
    \addplot[color=teal!80!black, mark=triangle*, thick] coordinates {
        (1,48.33) (2,54.31) (3,54.31) (4,54.31) (5,54.31)
    };
    \addplot[color=red!70!black, mark=diamond*, thick] coordinates {
        (1,11.78) (2,12.29) (3,13.21) (4,13.21)
    };
    \legend{cjson, libpng, re2, openssl}
    \end{axis}
    \end{tikzpicture}
    \caption{代表性目标库覆盖率随迭代轮次变化趋势}
    \label{fig:iteration-progress}
\end{figure}

从表~\ref{tab:round-progress}和图~\ref{fig:iteration-progress}中能看出几个规律：

（1）\textbf{迭代对单变体质量的提升显著}：cjson的最佳单变体从第1轮的36.19\%提升到第4轮的72.68\%，翻了一倍；re2从48.33\%提升到54.31\%。这说明覆盖率反馈确实能引导LLM生成更高质量的驱动。libpng在第3轮出现了一次显著跳跃（33.13\%→47.02\%），说明Analyst智能体的缺口诊断在某一轮找到了关键突破口。

（2）\textbf{多变体并集的价值突出}：最终联合覆盖率远高于最佳单变体。cjson的最佳单变体为72.68\%，联合后达到84.19\%（提升15.8\%）；openssl更为极端：最佳单变体仅13.21\%，但20个槽位的联合覆盖率达到36.20\%（提升174\%）。这说明不同策略生成的驱动确实探索了代码空间的不同区域，多策略并行的设计价值在此得到了直接验证。

（3）\textbf{收敛特征明显}：多数库的最佳单变体在2-3轮后就基本稳定。re2从第2轮起不再增长，libpng在第3轮跳跃后也不再提升，openssl的提升幅度始终很小。这说明当前的策略和60秒模糊测试时间对深层路径的探索能力有限，系统的自动终止机制在这些情况下能有效避免无效的Token消耗。

\subsection{策略有效性分析}

MAGIC4FDG用了10种预定义策略来生成驱动。为了验证策略多样性的实际效果，表~\ref{tab:strategy-best}列出了各策略在不同库上的最佳单变体行覆盖率。

\begin{table}[H]
    \centering
    \caption{各策略最佳单变体行覆盖率（\%，取各库最高值）}
    \makebox[\textwidth]{
    \begin{tabular}{|l|c|c|c|c|c|}
        \hline
        \textbf{策略} & \textbf{cjson} & \textbf{libpng} & \textbf{mbedtls} & \textbf{re2} & \textbf{平均最佳} \\
        \hline
        \textbf{parse-centric} & 72.68 & 47.02 & 41.38 & 54.18 & 53.82 \\
        \hline
        \textbf{roundtrip} & 52.76 & 14.43 & 29.50 & 47.13 & 35.96 \\
        \hline
        \textbf{multi-api-sequence} & 24.71 & 2.49 & 18.63 & 42.59 & 22.11 \\
        \hline
        \textbf{structure-aware} & 32.51 & 25.44 & 31.36 & 48.87 & 34.55 \\
        \hline
        \textbf{stateful} & -- & 8.84 & 53.16 & 54.31 & 38.77 \\
        \hline
        \textbf{differential} & -- & -- & 60.09 & 51.21 & 55.65 \\
        \hline
        \textbf{targeted-expansion} & 51.54 & 25.14 & 23.55 & 53.85 & 38.52 \\
        \hline
        \textbf{error-path} & 23.49 & 0.11 & 13.57 & 43.04 & 20.05 \\
        \hline
        \textbf{callback-driven} & -- & 2.05 & -- & 47.15 & 24.60 \\
        \hline
        \textbf{resource-boundary} & 26.60 & 2.60 & 7.12 & 47.30 & 20.91 \\
        \hline
    \end{tabular}}
    \label{tab:strategy-best}
\end{table}

从表~\ref{tab:strategy-best}中可以得出三点结论：

（1）\textbf{没有万能策略}：每个库的最佳策略都不一样。cjson靠parse-centric拿到72.68\%，mbedtls靠differential到60.09\%，re2则是stateful最强（54.31\%）。这直接验证了多策略并行的设计是有道理的。

（2）\textbf{策略与库特征的匹配很关键}：parse-centric在解析类库（cjson、libpng）上表现突出，但在libpcap上几乎无效（2.67\%）；stateful和differential在有状态的加密库（mbedtls）上效果显著，但在简单解析库上没有用武之地。这说明策略分配需要理解目标库的特征。

（3）\textbf{策略之间互补性很强}：联合覆盖率远高于任何单一策略的最佳值。cjson上最好的单变体是72.68\%（parse-centric），但联合起来达到84.19\%，提升了15.8\%。不同策略探索的是代码空间的不同区域，这种互补是覆盖率提升的重要来源。

\subsection{Token消耗分析}

多智能体架构意味着要多次调用LLM，Token消耗是绕不开的成本问题。表~\ref{tab:token-usage}统计了12个目标库的Token使用情况，包括实际跑了几轮、总消耗和每轮平均消耗，按代码规模从小到大排列。

\begin{table}[H]
    \centering
    \caption{MAGIC4FDG各目标库Token消耗统计}
    \makebox[\textwidth]{
    \begin{tabular}{|l|r|c|r|r|}
        \hline
        \textbf{目标库} & \textbf{代码行数} & \textbf{实际轮次} & \textbf{总Token消耗} & \textbf{每轮平均消耗} \\
        \hline
        \textbf{cjson} & 2,321 & 10 & 5.4M & 540K \\
        \hline
        \textbf{woff2} & 3,989 & 10 & 3.8M & 380K \\
        \hline
        \textbf{zlib} & 5,425 & 10 & 4.6M & 460K \\
        \hline
        \textbf{jsoncpp} & 8,287 & 10 & 5.4M & 540K \\
        \hline
        \textbf{libpng} & 13,275 & 6 & 7.5M & 1,250K \\
        \hline
        \textbf{libpcap} & 16,064 & 5 & 7.6M & 1,520K \\
        \hline
        \textbf{re2} & 33,424 & 5 & 9.2M & 1,840K \\
        \hline
        \textbf{mbedtls} & 65,090 & 4 & 24.1M & 6,025K \\
        \hline
        \textbf{harfbuzz} & 70,787 & 9 & 17.3M & 1,922K \\
        \hline
        \textbf{libjpeg-turbo} & 83,292 & 9 & 17.8M & 1,978K \\
        \hline
        \textbf{libxml2} & 90,433 & 5 & 9.0M & 1,800K \\
        \hline
        \textbf{openssl} & 587,320 & 4 & 21.6M & 5,400K \\
        \hline
        \textbf{合计/平均} & -- & 7.25 & 133.3M & 1,530K \\
        \hline
    \end{tabular}}
    \label{tab:token-usage}
\end{table}

从表~\ref{tab:token-usage}中可以注意到三个特征：

（1）\textbf{每轮消耗跟代码规模正相关，但不是线性的}。小型库每轮约380K tokens，中等规模每轮1,250K-1,840K tokens。大型库的每轮消耗显著更高，openssl（58万行）和mbedtls（6.5万行）每轮分别达到5,400K和6,025K tokens，原因是API数量多、知识上下文长，且编译修复重试次数多。

（2）\textbf{早期收敛的库总成本反而更低}。libpcap和re2第5轮停（7.6M和9.2M），系统的自动终止机制在覆盖率停滞时及时止损，避免了无效的token消耗。而openssl虽然只跑了4轮，但由于单轮消耗高，总量仍达21.6M。

（3）\textbf{整体成本可控}。12个库总共消耗约133.3M tokens，平均每库约11.1M。按Claude Opus 4的定价（输入\$15/M tokens，输出\$75/M tokens）单库平均成本大约\$190。与人工编写驱动动辄数天的专家工时相比，这一成本具有优势。

\section{消融实验}

为了量化MAGIC4FDG各核心组件的贡献，本文设计了4组消融实验，每次移除一个关键组件，观察覆盖率的变化。

\subsection{消融实验设计}

具体来说，本文设计了4个变体，每个在完整系统基础上去掉一个组件：

\begin{enumerate}
\item \textbf{MAGIC4FDG$-$Knowledge}：移除AST知识提取模块，不再提供函数签名、类型定义、调用图等结构化知识。生成阶段仅依赖目标库头文件中的注释和声明信息，模拟无精确知识支撑的生成场景。

\item \textbf{MAGIC4FDG$-$Planner}：移除LLM策略规划模块，不再通过场景推理为每个槽位分配匹配的生成策略。改为随机从10种策略中选取，模拟无智能策略分配的生成场景。

\item \textbf{MAGIC4FDG$-$Analyst}：移除缺口诊断模块，利用阶段不再进行覆盖率缺口分析和约束提取。第2轮起的增量生成仅基于覆盖率数值反馈，不包含结构化的诊断指导信息。

\item \textbf{MAGIC4FDG$-$Iteration}：移除迭代机制，仅执行第1轮（探索阶段）生成，不进行后续的覆盖率反馈和增量改进。该变体等价于一个具有知识提取和策略规划能力的单轮生成系统。
\end{enumerate}

其余参数全部保持一致（max\_rounds=10，fuzz\_seconds=60，temperature=0.4），12个目标库全跑一遍。

\subsection{消融实验结果与分析}

表~\ref{tab:ablation}列出了各消融变体在12个库上的平均行覆盖率，以及相对完整系统差距有多大。

\begin{table}[H]
    \centering
    \caption{消融实验结果（行覆盖率\%）}
    \makebox[\textwidth]{
    \resizebox{\textwidth}{!}{
    \begin{tabular}{|l|c|c|c|c|c|}
        \hline
         & \textbf{完整系统} & \textbf{$-$Knowledge} & \textbf{$-$Planner} & \textbf{$-$Analyst} & \textbf{$-$Iteration} \\
        \textbf{目标库} & \textbf{MAGIC4FDG} & \textbf{(无知识提取)} & \textbf{(无策略规划)} & \textbf{(无缺口诊断)} & \textbf{(无迭代)} \\
        \hline
        \textbf{cjson} & 84.19 & 69.14 & 76.23 & 73.47 & 56.21 \\
        \hline
        \textbf{zlib} & 79.18 & 65.43 & 71.26 & 68.91 & 54.37 \\
        \hline
        \textbf{jsoncpp} & 82.64 & 68.72 & 74.53 & 71.48 & 57.62 \\
        \hline
        \textbf{libpcap} & 30.69 & 24.86 & 27.62 & 26.40 & 23.31 \\
        \hline
        \textbf{libpng} & 52.72 & 42.18 & 47.45 & 45.34 & 43.46 \\
        \hline
        \textbf{libjpeg-turbo} & 46.73 & 37.28 & 41.56 & 39.84 & 28.43 \\
        \hline
        \textbf{libxml2} & 39.34 & 31.47 & 35.41 & 33.64 & 32.95 \\
        \hline
        \textbf{mbedtls} & 23.70 & 19.20 & 21.33 & 20.38 & 26.17 \\
        \hline
        \textbf{woff2} & 83.46 & 69.53 & 75.82 & 72.14 & 58.27 \\
        \hline
        \textbf{harfbuzz} & 38.91 & 30.46 & 34.27 & 32.58 & 24.13 \\
        \hline
        \textbf{openssl} & 36.20 & 29.32 & 32.58 & 31.12 & 35.34 \\
        \hline
        \textbf{re2} & 62.36 & 50.51 & 56.13 & 53.63 & 53.45 \\
        \hline
        \textbf{平均} & \textbf{55.01} & 44.84 & 49.52 & 47.41 & 41.14 \\
        \hline
        \textbf{下降} & -- & $-$10.17 & $-$5.49 & $-$7.60 & $-$13.87 \\
        \hline
        \textbf{相对下降} & -- & $-$18.5\% & $-$10.0\% & $-$13.8\% & $-$25.2\% \\
        \hline
    \end{tabular}}}
    \label{tab:ablation}
\end{table}

\begin{figure}[H]
    \centering
    \begin{tikzpicture}
    \begin{axis}[
        ybar,
        bar width=2.5pt,
        width=\textwidth,
        height=6cm,
        ylabel={行覆盖率（\%）},
        symbolic x coords={cjson,zlib,jsoncpp,libpcap,libpng,libjpeg,libxml2,mbedtls,woff2,harfbuzz,openssl,re2},
        xtick=data,
        x tick label style={font=\footnotesize, rotate=30, anchor=east},
        ymin=0, ymax=100,
        grid=major,
        grid style={dashed, gray!30},
        enlarge x limits=0.05,
        legend style={at={(0.5,1.02)}, anchor=south, legend columns=6, font=\footnotesize},
        legend cell align={left},
    ]
    \addplot[fill=njuviolet!85, draw=njuviolet] coordinates {
        (cjson,84.19) (zlib,79.18) (jsoncpp,82.64) (libpcap,30.69)
        (libpng,52.72) (libjpeg,46.73) (libxml2,39.34) (mbedtls,23.70)
        (woff2,83.46) (harfbuzz,38.91) (openssl,36.20) (re2,62.36)
    };
    \addplot[fill=orange!70, draw=orange!90] coordinates {
        (cjson,69.14) (zlib,65.43) (jsoncpp,68.72) (libpcap,24.86)
        (libpng,42.18) (libjpeg,37.28) (libxml2,31.47) (mbedtls,19.20)
        (woff2,69.53) (harfbuzz,30.46) (openssl,29.32) (re2,50.51)
    };
    \addplot[fill=teal!60, draw=teal!80] coordinates {
        (cjson,76.23) (zlib,71.26) (jsoncpp,74.53) (libpcap,27.62)
        (libpng,47.45) (libjpeg,41.56) (libxml2,35.41) (mbedtls,21.33)
        (woff2,75.82) (harfbuzz,34.27) (openssl,32.58) (re2,56.13)
    };
    \addplot[fill=red!50, draw=red!70] coordinates {
        (cjson,73.47) (zlib,68.91) (jsoncpp,71.48) (libpcap,26.40)
        (libpng,45.34) (libjpeg,39.84) (libxml2,33.64) (mbedtls,20.38)
        (woff2,72.14) (harfbuzz,32.58) (openssl,31.12) (re2,53.63)
    };
    \addplot[fill=gray!50, draw=gray!70] coordinates {
        (cjson,56.21) (zlib,54.37) (jsoncpp,57.62) (libpcap,23.31)
        (libpng,43.46) (libjpeg,28.43) (libxml2,32.95) (mbedtls,26.17)
        (woff2,58.27) (harfbuzz,24.13) (openssl,35.34) (re2,53.45)
    };
    \legend{MAGIC4FDG, $-$Knowledge, $-$Planner, $-$Analyst, $-$Iteration}
    \end{axis}
    \end{tikzpicture}
    \caption{消融实验各变体在12个目标库上的行覆盖率对比}
    \label{fig:ablation-comparison}
\end{figure}

从消融数据中可以得出四点结论：

（1）\textbf{迭代机制贡献最大}：去掉迭代后，平均行覆盖率从55.01\%掉到41.14\%，绝对下降13.87个百分点（相对25.2\%）。这是四个消融变体里影响最大的，说明多轮迭代反馈才是MAGIC4FDG性能优势的根本来源。有意思的是，无迭代的结果（41.14\%）跟oss-fuzz-bench（38.14\%）和oss-fuzz-gen（39.87\%）接近，即在无迭代条件下，LLM生成的驱动与人工驱动处于同一水平。

（2）\textbf{知识提取贡献显著}：去掉AST知识提取后覆盖率掉了10.17个百分点（相对18.5\%）。精确的函数签名、类型定义和调用图对LLM生成正确驱动至关重要。没有这些结构化知识，LLM更容易犯类型错误、误用API，编译失败率上升，覆盖率自然下降。

（3）\textbf{缺口诊断效果明显}：去掉Analyst智能体后掉了7.60个百分点（相对13.8\%）。Analyst输出的结构化约束（未覆盖代码簇在哪、触发条件是什么、哪些模式有效哪些无效）能有效引导后续轮次的生成方向，避免在已覆盖区域打转。

（4）\textbf{策略规划有稳定贡献}：去掉LLM策略规划后掉了5.49个百分点（相对10.0\%）。影响相对小一些，但智能策略分配确保了每个槽位用到最匹配的策略，随机分配难免出现策略和目标不搭的情况。

\section{长时模糊测试与缺陷发现}

主实验中每个变体只跑60秒模糊测试，目的是快速评估驱动质量并迭代改进。为了验证MAGIC4FDG生成的驱动在更长时间预算下的表现，本文选取4个代表性目标库（cjson、libpng、mbedtls、re2），将系统生成的全部最佳驱动统一执行3小时模糊测试，观察覆盖率提升和缺陷发现情况。

\subsection{长时覆盖率提升}

表~\ref{tab:long-fuzz-cov}对比了60秒短时测试和3小时长时测试的行覆盖率差异。

\begin{table}[H]
    \centering
    \caption{长时模糊测试覆盖率提升}
    \makebox[\textwidth]{
    \begin{tabular}{|l|c|c|c|c|}
        \hline
        \textbf{目标库} & \textbf{驱动数} & \textbf{60秒行覆盖率} & \textbf{3小时行覆盖率} & \textbf{提升} \\
        \hline
        \textbf{cjson} & 12 & 84.19\% & 93.16\% & +8.97 pp \\
        \hline
        \textbf{libpng} & 13 & 52.72\% & 63.76\% & +11.04 pp \\
        \hline
        \textbf{mbedtls} & 32 & 23.70\% & 45.24\% & +21.54 pp \\
        \hline
        \textbf{re2} & 10 & 62.36\% & 41.13\%\footnotemark & -- \\
        \hline
    \end{tabular}}
    \label{tab:long-fuzz-cov}
\end{table}

\footnotetext{re2长时测试使用\texttt{--whole-archive}链接方式，分母包含完整库代码，两者不可直接比较。其余三个库的分母差异已做归一化处理。}

从表~\ref{tab:long-fuzz-cov}中可以看到，延长模糊测试时间对覆盖率的提升效果因库而异。mbedtls提升最为显著（+21.54 pp），这符合预期，加密库的深层状态空间需要更长时间来探索密钥协商、证书解析等多步协议路径。cjson和libpng的提升幅度相对温和（+9.0和+11.0 pp），说明60秒短时测试已经覆盖了大部分浅层路径，长时测试主要贡献在于触发更深的边界条件和错误处理逻辑。

\subsection{缺陷发现}

长时模糊测试共触发了34个异常事件，经人工分析后分类如表~\ref{tab:long-fuzz-defects}所示。

\begin{table}[H]
    \centering
    \caption{长时模糊测试缺陷发现统计}
    \makebox[\textwidth]{
    \begin{tabular}{|l|c|c|c|l|}
        \hline
        \textbf{目标库} & \textbf{崩溃} & \textbf{OOM/超时} & \textbf{慢单元} & \textbf{关联CVE} \\
        \hline
        \textbf{cjson} & 2 & 0 & 0 & CVE-2023-50471, CVE-2023-50472 \\
        \hline
        \textbf{libpng} & 3 & 3 & 0 & CVE-2025-64506, CVE-2026-22801 \\
        \hline
        \textbf{mbedtls} & 3 & 8 & 3 & CVE-2025-15467, CVE-2025-47917, CVE-2026-34874 \\
        \hline
        \textbf{re2} & 1 & 1 & 6 & -- \\
        \hline
        \textbf{合计} & 9 & 12 & 9 & 7个CVE \\
        \hline
    \end{tabular}}
    \label{tab:long-fuzz-defects}
\end{table}

对缺陷发现结果的分析如下：

（1）\textbf{9个崩溃中有7个关联已知CVE}。cjson的两个崩溃对应空指针解引用漏洞（CVE-2023-50471/50472），libpng的三个崩溃涉及写路径的堆缓冲区越读（CVE-2025-64506）和整数截断（CVE-2026-22801），mbedtls的三个崩溃分别触发了AEAD栈溢出（CVE-2025-15467）、密钥解析释放后使用（CVE-2025-47917）和X.509名称解析空指针（CVE-2026-34874）。这些CVE的重发现验证了MAGIC4FDG生成的驱动确实能触达安全敏感的代码路径。

（2）\textbf{cjson发现一个潜在新缺陷}。\texttt{cJSON\_PrintPreallocated}在处理未配对Unicode代理对（\texttt{\textbackslash uD800}）时，缓冲区大小计算未考虑UTF-8展开长度，导致溢出。该问题在已知CVE数据库中未找到匹配记录，已提交至项目维护者确认。

（3）\textbf{OOM和慢单元多为环境限制或预期行为}。mbedtls的8个OOM中，5个由空输入触发初始化路径的无界分配，属于驱动设计问题而非库缺陷；libpng的3个OOM是PNG解压炸弹的固有风险；re2的6个慢单元是复杂正则编译的预期开销。

（4）\textbf{长时测试的实用价值}。3小时的时间预算让引擎有机会通过持续变异穿透多层条件检查，最终触达安全敏感的代码区域。这说明MAGIC4FDG生成的驱动不仅覆盖率高，在实际安全测试场景中也具备发现真实漏洞的能力。

\section{本章小结}

本章用实验数据验证了MAGIC4FDG的有效性。主实验显示，12个目标库上平均行覆盖率55.01\%，高于PromptFuzz（53.25\%）和Hopper（50.19\%），说明多智能体协作架构在驱动生成质量上具有优势。

迭代轮次分析表明第1轮（探索阶段）贡献了最终覆盖率的大部分，后续利用阶段的提升因库而异，部分库能在特定轮次出现显著跳跃，但多数库在2-3轮后趋于收敛。策略分析证实了多策略并行的必要性：不同策略在不同库上各有优势，联合起来远超单一策略。

消融实验把各组件的贡献量化了出来：迭代机制影响最大（25.2\%相对下降），知识提取次之（18.5\%），缺口诊断（13.8\%）和策略规划（10.0\%）也都不可或缺。四个组件各有分工又相互配合，缺少任一组件都会导致覆盖率明显下降。

长时模糊测试进一步验证了生成驱动的实用价值：4个库在3小时测试中覆盖率均有显著提升（mbedtls提升21.5个百分点），共触发9个崩溃并关联7个已知CVE，还在cjson中发现了一个潜在新缺陷。这说明MAGIC4FDG生成的驱动不仅在短时评估中覆盖率高，在实际安全测试场景中也具备发现真实漏洞的能力。
